% =====================================================================
% APÊNDICE B — PARÂMETROS FÍSICOS DO PERSEGUIDOR (CHASER)
% Fonte: Parametros_chaser.m
% Padrão:
%   - SI estrito (m, kg, s, kg·m^2)
%   - 2 casas decimais quando aplicável; inércias em notação científica
%   - tabelas sempre com [H]
% Requer: \usepackage{float}
% =====================================================================

% \section{Parâmetros físicos do perseguidor}
\label{ap:params_chaser}

As tabelas deste apêndice consolidam os parâmetros físicos do perseguidor.

% ------------------------------------------------------------
\subsection{Geometria, massa e centros de massa (SI)}

\begin{table}[H]
\centering
\caption{Base geral: massa, dimensões e centro de massa (SI).}
\label{tab:chaser_base_SI}
\begin{tabular}{lcc}
\hline
\textbf{Grandeza} & \textbf{Símbolo} & \textbf{Valor} \\
\hline
Massa & $m$ (kg) & 3,65 \\
Comprimento & $L$ (m) & 0,10 \\
Raio & $R$ (m) & 0,25 \\
Diâmetro & $D$ (m) & 0,50 \\
Centro de massa & $\vec{r}_{cm}$ (m) & $[\,0,00,\ 0,00,\ 0,00\,]^T$ \\
\hline
\end{tabular}
\end{table}

\begin{table}[H]
\centering
\caption{Juntas: massa, dimensões e centros de massa (SI).}
\label{tab:chaser_juntas_SI}
\begin{tabular}{lcccc}
\hline
\textbf{Junta} & \textbf{Massa (kg)} & \textbf{$L$ (m)} & \textbf{$R$ (m)} & \textbf{$\vec{r}_{cm}$ (m)} \\
\hline
J1 & 0,07 & 0,05 & 0,05 & $[\,0,00,\ -0,13,\ 0,08\,]^T$ \\
J2 & 0,07 & 0,05 & 0,05 & $[\,0,00,\ -0,13,\ 0,08\,]^T$ \\
J3 & 0,07 & 0,05 & 0,05 & $[\,0,00,\ 0,23,\ 0,23\,]^T$ \\
\hline
\end{tabular}
\end{table}

\begin{table}[H]
\centering
\caption{Elos: massa, dimensões e centros de massa (SI).}
\label{tab:chaser_elos_SI}
\begin{tabular}{lcccc}
\hline
\textbf{Elo} & \textbf{Massa (kg)} & \textbf{$L$ (m)} & \textbf{$R$ (m)} & \textbf{$\vec{r}_{cm}$ (m)} \\
\hline
E1 & 0,12 & 0,10 & 0,03 & $[\,0,00,\ -0,13,\ 0,15\,]^T$ \\
E2 & 0,32 & 0,30 & 0,03 & $[\,0,00,\ 0,05,\ 0,23\,]^T$ \\
E3 & 0,32 & 0,30 & 0,03 & $[\,0,00,\ 0,40,\ 0,23\,]^T$ \\
\hline
\end{tabular}
\end{table}

\begin{table}[H]
\centering
\caption{Ferramenta: massa, dimensões e centro de massa (SI).}
\label{tab:chaser_ferramenta_SI}
\begin{tabular}{lcc}
\hline
\textbf{Grandeza} & \textbf{Símbolo} & \textbf{Valor} \\
\hline
Massa & $m$ (kg) & 0,04 \\
Comprimento & $L$ (m) & 0,15 \\
Raio & $R$ (m) & 0,01 \\
Diâmetro & $D$ (m) & 0,02 \\
Centro de massa & $\vec{r}_{cm}$ (m) & $[\,0,00,\ 0,63,\ 0,23\,]^T$ \\
\hline
\end{tabular}
\end{table}

\begin{table}[H]
\centering
\caption{Conjuntos combinados: massas e centros de massa (SI).}
\label{tab:chaser_combinados_SI}
\begin{tabular}{lcc}
\hline
\textbf{Conjunto} & \textbf{Massa (kg)} & \textbf{$\vec{r}_{cm}$ (m)} \\
\hline
J1 + E1 & 0,19 & $[\,0,00,\ -0,13,\ 0,12\,]^T$ \\
J2 + E2 & 0,40 & $[\,0,00,\ 0,02,\ 0,23\,]^T$ \\
J3 + E3 + Ferramenta & 0,44 & $[\,0,00,\ 0,39,\ 0,23\,]^T$ \\
\hline
\end{tabular}
\end{table}

\begin{table}[H]
\centering
\caption{Chaser (home, modo 0): centro de massa (SI).}
\label{tab:chaser_home_com_SI}
\begin{tabular}{lc}
\hline
\textbf{Grandeza} & \textbf{Valor} \\
\hline
$\vec{r}_{cm}$ (m) & $[\,0,00,\ 0,03,\ 0,05\,]^T$ \\
\hline
\end{tabular}
\end{table}

\begin{table}[H]
\centering
\caption{Espessuras adotadas para cascas (SI).}
\label{tab:chaser_espessuras_SI}
\begin{tabular}{lc}
\hline
\textbf{Elemento} & \textbf{Espessura (m)} \\
\hline
Parede da base & 0,00 \\
Parede dos elos & 0,00 \\
Parede da ferramenta & 0,00 \\
\hline
\end{tabular}
\end{table}

% ------------------------------------------------------------
\subsection{Tensores de inércia (SI)}

% OBS: Para caber na página, estas tabelas foram quebradas por componente
% (uma linha por tabela). Mantém leitura direta e evita largura excessiva.

\begin{table}[H]
\centering
\caption{Tensor de inércia da base geral (kg$\cdot$m$^2$).}
\label{tab:chaser_I_base_SI}
\begin{tabular}{lccc}
\hline
$I_{xx}$ & $I_{yy}$ & $I_{zz}$ & Termos cruzados $(I_{xy}, I_{xz}, I_{yz})$ \\
\hline
$7,95\times 10^{-2}$ & $7,95\times 10^{-2}$ & $1,45\times 10^{-1}$ &
$(6,60\times 10^{-16},\ 1,64\times 10^{-12},\ -1,26\times 10^{-16})$ \\
\hline
\end{tabular}
\end{table}

\begin{table}[H]
\centering
\caption{Tensores de inércia dos elos (kg$\cdot$m$^2$): parte 1/2.}
\label{tab:chaser_I_elos_1_SI}
\begin{tabular}{lcccc}
\hline
\textbf{Componente} & $I_{xx}$ & $I_{yy}$ & $I_{zz}$ & $(I_{xy}, I_{xz}, I_{yz})$ \\
\hline
Elo 1 &
$1,66\times 10^{-4}$ &
$1,66\times 10^{-4}$ &
$6,24\times 10^{-5}$ &
$(0,\ 5,74\times 10^{-14},\ 0)$ \\
Elo 2 &
$2,83\times 10^{-3}$ &
$1,76\times 10^{-4}$ &
$2,83\times 10^{-3}$ &
$(0,\ 1,88\times 10^{-17},\ 0)$ \\
\hline
\end{tabular}
\end{table}

\begin{table}[H]
\centering
\caption{Tensores de inércia dos elos e ferramenta (kg$\cdot$m$^2$): parte 2/2.}
\label{tab:chaser_I_elos_2_SI}
\begin{tabular}{lcccc}
\hline
\textbf{Componente} & $I_{xx}$ & $I_{yy}$ & $I_{zz}$ & $(I_{xy}, I_{xz}, I_{yz})$ \\
\hline
Elo 3 &
$2,83\times 10^{-3}$ &
$1,76\times 10^{-4}$ &
$2,83\times 10^{-3}$ &
$(1,25\times 10^{-17},\ 0,\ 0)$ \\
Ferramenta &
$8,11\times 10^{-5}$ &
$1,63\times 10^{-6}$ &
$8,11\times 10^{-5}$ &
$(0,\ 0,\ 0)$ \\
\hline
\end{tabular}
\end{table}

% =====================================================================
% AJUSTE — QUEBRA DE TABELAS (SEM ESTOURAR LARGURA)
% Sempre [H]. Tensores dos conjuntos combinados divididos por conjunto
% =====================================================================

\subsection{Tensores de inércia dos conjuntos combinados (SI)}
\label{sec:apB_inercias_combinadas}

% --------- B.12: J1 + E1 ---------
\begin{table}[H]
\centering
\caption{Tensor de inércia do conjunto J1 + E1 (kg$\cdot$m$^2$).}
\label{tab:B12_I_J1E1}
\begin{tabular}{lc}
\hline
\textbf{Elemento} & \textbf{Valor} \\
\hline
$I_{xx}$ & $4,60\times 10^{-4}$ \\
$I_{yy}$ & $4,60\times 10^{-4}$ \\
$I_{zz}$ & $9,63\times 10^{-5}$ \\
$I_{xy}$ & $0$ \\
$I_{xz}$ & $-1,04\times 10^{-14}$ \\
$I_{yz}$ & $0$ \\
\hline
\end{tabular}
\end{table}

% --------- B.13: J2 + E2 ---------
\begin{table}[H]
\centering
\caption{Tensor de inércia do conjunto J2 + E2 (kg$\cdot$m$^2$).}
\label{tab:B13_I_J2E2}
\begin{tabular}{lc}
\hline
\textbf{Elemento} & \textbf{Valor} \\
\hline
$I_{xx}$ & $4,66\times 10^{-3}$ \\
$I_{yy}$ & $2,16\times 10^{-4}$ \\
$I_{zz}$ & $4,67\times 10^{-3}$ \\
$I_{xy}$ & $-4,44\times 10^{-13}$ \\
$I_{xz}$ & $3,12\times 10^{-17}$ \\
$I_{yz}$ & $-1,49\times 10^{-17}$ \\
\hline
\end{tabular}
\end{table}

% --------- B.14: J3 + E3 + Ferramenta ---------
\begin{table}[H]
\centering
\caption{Tensor de inércia do conjunto J3 + E3 + Ferramenta (kg$\cdot$m$^2$).}
\label{tab:B14_I_J3E3F}
\begin{tabular}{lc}
\hline
\textbf{Elemento} & \textbf{Valor} \\
\hline
$I_{xx}$ & $7,18\times 10^{-3}$ \\
$I_{yy}$ & $2,18\times 10^{-4}$ \\
$I_{zz}$ & $7,19\times 10^{-3}$ \\
$I_{xy}$ & $-3,78\times 10^{-13}$ \\
$I_{xz}$ & $0$ \\
$I_{yz}$ & $-7,60\times 10^{-14}$ \\
\hline
\end{tabular}
\end{table}


\begin{table}[H]
\centering
\caption{Tensor de inércia do chaser no modo home (kg$\cdot$m$^2$).}
\label{tab:chaser_I_home_SI}
\begin{tabular}{lcccc}
\hline
$I_{xx}$ & $I_{yy}$ & $I_{zz}$ & $(I_{xy}, I_{xz}, I_{yz})$ \\
\hline
$2,11\times 10^{-1}$ &
$1,34\times 10^{-1}$ &
$2,22\times 10^{-1}$ &
$(-7,64\times 10^{-16},\ -5,86\times 10^{-16},\ -3,62\times 10^{-2})$ \\
\hline
\end{tabular}
\end{table}

% ------------------------------------------------------------
\subsection{Massas agregadas}

\begin{table}[H]
\centering
\caption{Massas agregadas do perseguidor (SI).}
\label{tab:chaser_massas_agregadas_SI}
\begin{tabular}{lc}
\hline
\textbf{Grandeza} & \textbf{Valor (kg)} \\
\hline
Massa total do chaser (base + MR) & 4,68 \\
Massa total dos elos (E1 + E2 + E3) & 0,77 \\
Massa total das juntas (J1 + J2 + J3) & 0,22 \\
Massa total (elos + juntas) & 0,99 \\
Massa total do MR (elos + juntas + ferramenta) & 1,03 \\
\hline
\end{tabular}
\end{table}

% ------------------------------------------------------------
\subsection{Limites das juntas do manipulador robótico}

\begin{table}[H]
\centering
\caption{Limites de posição das juntas (SI).}
\label{tab:chaser_lim_pos_SI}
\begin{tabular}{lcc}
\hline
\textbf{Junta} & \textbf{Posição (rad)} & \textbf{Posição (deg)} \\
\hline
J1 & $[-6,28,\ 6,28]$ & $[-360,00,\ 360,00]$ \\
J2 & $[-1,57,\ 1,57]$ & $[-90,00,\ 90,00]$ \\
J3 & $[-2,27,\ 2,27]$ & $[-130,00,\ 130,00]$ \\
\hline
\end{tabular}
\end{table}

\begin{table}[H]
\centering
\caption{Limites de velocidade angular das juntas (SI).}
\label{tab:chaser_lim_vel_SI}
\begin{tabular}{lc}
\hline
\textbf{Junta} & \textbf{$|\dot{q}|_{\max}$ (rad/s)} \\
\hline
J1 & 0,31 \\
J2 & 0,31 \\
J3 & 0,31 \\
\hline
\end{tabular}
\end{table}

\begin{table}[H]
\centering
\caption{Limites de torque das juntas (SI).}
\label{tab:chaser_lim_tau_SI}
\begin{tabular}{lcc}
\hline
\textbf{Junta} & \textbf{$\tau$ nominal (Nm)} & \textbf{$\tau$ máximo (Nm)} \\
\hline
J1 & $[-4,00,\ 4,00]$ & $[-4,20,\ 4,20]$ \\
J2 & $[-4,00,\ 4,00]$ & $[-4,20,\ 4,20]$ \\
J3 & $[-4,00,\ 4,00]$ & $[-4,20,\ 4,20]$ \\
\hline
\end{tabular}
\end{table}
